\section*{Chapitre \ref{ch:b1:enzymes} --- Enzymes et catalyse biochimique}

\begin{solution}{exo:b1:enzymes:1}
Elle abaisse l'\omterm{prop:b1:enzymes:whatitdoes}{énergie d'activation} (la hauteur de l'\omterm{prop:b1:enzymes:whatitdoes}{état de
transition}), et donc la vitesse, dans les deux sens~; elle laisse
inchangés $\Delta G$, la constante d'équilibre et la position de
l'équilibre. Sur le profil, le sommet est abaissé, les deux extrémités
non.
\end{solution}

\begin{solution}{exo:b1:enzymes:2}
$K_m$~: concentration de \omterm{def:b1:enzymes:enzyme}{substrat} à la moitié de $V_{\max}$
(\unit{mol/L}). $V_{\max}$~: vitesse à \omterm{def:b1:enzymes:enzyme}{substrat} saturant (\unit{mol/s},
ou \unit{\micro
mol/min}). $k_{\text{cat}}$~: molécules de \omterm{def:b1:enzymes:enzyme}{substrat} converties par
seconde et par \omterm{def:b1:enzymes:enzyme}{enzyme} à saturation (\unit{s^{-1}}).
$k_{\text{cat}}/K_m$~: l'efficacité, constante de vitesse du second
ordre à faible \omterm{def:b1:enzymes:enzyme}{substrat} (\unit{L.mol^{-1}.s^{-1}}).
\end{solution}

\begin{solution}{exo:b1:enzymes:3}
Sans \omterm{def:b1:enzymes:inhibitors}{inhibiteur}~: ordonnée à l'origine 1, donc $V_{\max} = 1$~;
abscisse $x$ à l'origine $-0.5$, donc $K_m = 2$. \omterm{def:b1:enzymes:inhibitors}{Compétitif}~:
$V_{\max} = 1$, abscisse $x$ $-0.25$, $K_m = 4$. Non \omterm{def:b1:enzymes:inhibitors}{compétitif}~:
ordonnée 2, $V_{\max} = 0.5$~; abscisse $x$ $-0.5$, $K_m = 2$.
\end{solution}

\begin{solution}{exo:b1:enzymes:4}
Régulation \omterm{def:b1:proteins:allostery}{allostérique} (\omterm{prop:b1:enzymes:control}{rétro-inhibition}), quelques millisecondes~;
\omterm{prop:b1:enzymes:control}{modification covalente} (phosphorylation), de la seconde à la minute~;
activation protéolytique d'un zymogène, une fois, à la demande~;
contrôle de la quantité par l'expression du gène et la dégradation, de
la minute à l'heure.
\end{solution}

\begin{solution}{exo:b1:enzymes:5}
$v_0 = 50[S]/(0.1 + [S])$~: 8.3, 25, 45.5 et \qty{49.5}{\micro mol/min}.
$0.9 = [S]/(0.1 + [S])$ donne $[S] = 0.9\,K_m/0.1 = \qty{0.9}{mmol/L}$.
\end{solution}

\begin{solution}{exo:b1:enzymes:6}
$e^{20\,000/2580} = e^{7.75} = 2300$. Pour $10^{10}$~: $\Delta E_a =
RT\ln 10^{10} = 2.58\times 23.0 = \qty{59}{kJ/mol}$.
\end{solution}

\begin{solution}{exo:b1:enzymes:7}
$V_{\max} = \qty{0.12}{mol/L}$ par minute dans \qty{1}{mL}~:
\qty{1.2e-4}{mol/min} $= \qty{2.0e-6}{mol/s}$ de produit~; \omterm{def:b1:enzymes:enzyme}{enzyme}
\qty{2e-12}{mol}~: $k_{\text{cat}} = 2\times 10^{-6}/2\times 10^{-12} =
\qty{1e6}{s^{-1}}$. Efficacité $10^6/0.012 = \qty{8e7}{L.mol^{-1}.s^{-1}}$,
à un facteur dix près de la limite de diffusion~: presque chaque
rencontre avec une molécule de \omterm{def:b1:enzymes:enzyme}{substrat} est productive.
\end{solution}

\begin{solution}{exo:b1:enzymes:8}
$1 + 2/K_i = 2$~: $K_i = \qty{2}{mmol/L}$. Pour \qty{90}{\%}~:
$[S] = 9K_m^{\text{app}}$~: \qty{9}{mmol/L} sans l'\omterm{def:b1:enzymes:inhibitors}{inhibiteur},
\qty{18}{mmol/L} avec.
\end{solution}

\begin{solution}{exo:b1:enzymes:9}
Inhiber la première étape engagée arrête toute la voie d'un coup et ne
gaspille aucun intermédiaire, qui s'accumulerait sinon~; le produit
final est le signal qui compte --- son abondance est exactement
l'information que la voie n'est plus nécessaire --- alors que le taux
d'un intermédiaire ne dit rien de la demande.
\end{solution}

\begin{solution}{exo:b1:enzymes:10}
Une protéase active dans la \omterm{def:b1:cell-unit-of-life:cell}{cellule} qui la fabrique digérerait cette
\omterm{def:b1:cell-unit-of-life:cell}{cellule}~; fabriquée en zymogène inactif, elle est inoffensive jusqu'à
ce que l'entéropeptidase, \omterm{def:b1:enzymes:enzyme}{enzyme} du revêtement intestinal, coupe le
trypsinogène en trypsine dans le duodénum, là où l'on veut digérer~; la
trypsine active ensuite les autres zymogènes (une cascade). Comme une
trace de trypsine peut déclencher la cascade prématurément, le pancréas
emballe dans ses granules un \omterm{def:b1:enzymes:inhibitors}{inhibiteur} spécifique de la trypsine, par
sécurité~; la défaillance de ce dispositif est la pancréatite.
\end{solution}

\begin{solution}{exo:b1:enzymes:11}
Sans~: $3^3/(3^3 + 3^3) = 0.50$. Avec~: $27/(166 + 27) = 0.14$~: la
vitesse tombe à \qty{28}{\%} de sa valeur. Une \omterm{def:b1:enzymes:enzyme}{enzyme} michaélienne dont
le $K_m$ passe de 3 à 5.5~: $3/6 = 0.50$ devient $3/8.5 = 0.35$, soit
\qty{70}{\%}. La \omterm{def:b1:proteins:allostery}{coopérativité} rend le même décalage de courbe plus de
deux fois plus efficace au voisinage du point de fonctionnement~: un
interrupteur plutôt qu'un variateur.
\end{solution}

\begin{solution}{exo:b1:enzymes:12}
Un simple catalyseur (une surface de platine, un acide) accélère
indistinctement de nombreuses réactions~; une \omterm{def:b1:enzymes:enzyme}{enzyme} accélère une
réaction sur un \omterm{def:b1:enzymes:enzyme}{substrat} --- la spécificité de son \omterm{def:b1:enzymes:enzyme}{site actif} est déjà
une information sur ce que la \omterm{def:b1:cell-unit-of-life:cell}{cellule} veut voir faire. La régulation
ajoute une seconde couche~: sites \omterm{def:b1:proteins:allostery}{allostériques}, sites de
phosphorylation et peptides des zymogènes disent à l'\omterm{def:b1:enzymes:enzyme}{enzyme} quand et où
agir, en réponse à l'état de la \omterm{def:b1:cell-unit-of-life:cell}{cellule}. Le coût est une grosse \omterm{def:b1:proteins:peptide}{protéine}
(des centaines de résidus pour un site d'une douzaine), un
renouvellement plus lent que celui des meilleurs catalyseurs
inorganiques, et la machinerie pour la fabriquer, la modifier et la
détruire~; le bénéfice est une chimie qui ne tourne que là et quand elle
est utile, et c'est cela, un métabolisme.
\end{solution}

\begin{solution}{pb:b1:enzymes:1}
\textbf{1.} $1/[S]$~: 2, 1, 0.5, 0.2, 0.1, 0.05. $1/v_0$~: 5.0, 3.0, 2.0,
1.4, 1.2, 1.1.
\textbf{2.} Chaque pas de $1/[S]$ change $1/v_0$ proportionnellement~:
pente $(5.0 - 1.1)/(2 - 0.05) = 2.0$~; ordonnée à l'origine $1.0$.
\textbf{3.} $V_{\max} = 1/1.0 = \qty{1.0}{\micro mol/min}$~; $K_m =
\text{pente}\times V_{\max} = \qty{2.0}{mmol/L}$.
\textbf{4.} $1.0\times 5/(2 + 5) = 0.714$~: la valeur mesurée.
\textbf{5.} $\qty{10e-9}{mol/L}\times\qty{1e-3}{L} = \qty{1e-11}{mol}$~;
$V_{\max} = \qty{1e-6}{mol/min} = \qty{1.67e-8}{mol/s}$~;
$k_{\text{cat}} = 1.67\times 10^{-8}/10^{-11} = \qty{1670}{s^{-1}}$.
\textbf{6.} $1670/(2\times 10^{-3}) = \qty{8.3e5}{L.mol^{-1}.s^{-1}}$~:
mille fois sous la limite de diffusion~; une collision sur mille
seulement est productive.
\textbf{7.} $0.95 = [S]/(2 + [S])$~: $[S] = 19\,K_m = \qty{38}{mmol/L}$.
\textbf{8.} $1/v_0$~: 9.0, 5.0, 3.0, 1.8, 1.4, 1.2. Pente $(9.0 -
1.2)/1.95 = 4.0$~; ordonnée à l'origine $1.0$.
\textbf{9.} $V_{\max}^{\text{app}} = \qty{1.0}{\micro mol/min}$
(inchangé)~; $K_m^{\text{app}} = 4.0\times 1.0 = \qty{4.0}{mmol/L}$.
\textbf{10.} $K_m$ doublé, $V_{\max}$ inchangé~: \omterm{def:b1:enzymes:inhibitors}{compétitif}.
\textbf{11.} $4 = 2(1 + 1/K_i)$~: $K_i = \qty{1.0}{mmol/L}$.
\textbf{12.} Avec l'\omterm{def:b1:enzymes:inhibitors}{inhibiteur}, $v = V_{\max}[S]/(K_m(1 + [I]/K_i) +
[S])$~; diviser la vitesse par deux exige $K_m(1 + [I]/K_i) + [S] = 2(K_m +
[S])$, soit $[I] = K_i(K_m + [S])/K_m$~: à \qty{2}{mmol/L}, $[I] =
1\times 4/2 = \qty{2}{mmol/L}$~; à \qty{20}{mmol/L}, $[I] = 22/2 =
\qty{11}{mmol/L}$. Plus il y a de \omterm{def:b1:enzymes:enzyme}{substrat}, plus il faut d'\omterm{def:b1:enzymes:inhibitors}{inhibiteur}.
\textbf{13.} Une molécule ressemblant au \omterm{def:b1:enzymes:enzyme}{substrat} (ou à son \omterm{prop:b1:enzymes:whatitdoes}{état de
transition}) --- un analogue d'ester non hydrolysable --- qui se fixe
dans le \omterm{def:b1:enzymes:enzyme}{site actif}.
\textbf{14.} $1/v_0$~: 10.0, 6.0, 4.0, 2.8, 2.4, 2.2~; pente $4.0$,
ordonnée à l'origine $2.0$~: $V_{\max}^{\text{app}} = \qty{0.5}{\micro mol/min}$,
$K_m^{\text{app}} = 4.0\times 0.5 = \qty{2.0}{mmol/L}$.
\textbf{15.} $V_{\max}$ divisé par deux, $K_m$ inchangé~: non
\omterm{def:b1:enzymes:inhibitors}{compétitif}~; $2 = 1 + 1/K_i$~: $K_i = \qty{1.0}{mmol/L}$.
\textbf{16.} J se fixe sur un site autre que le \omterm{def:b1:enzymes:enzyme}{site actif}, sur l'\omterm{def:b1:enzymes:enzyme}{enzyme}
libre comme sur $ES$, avec la même affinité~: le \omterm{def:b1:enzymes:enzyme}{substrat} ne lui fait
pas concurrence, et à tout $[S]$ la moitié des molécules d'\omterm{def:b1:enzymes:enzyme}{enzyme} est
inactivée.
\textbf{17.} Avec J la vitesse double (la moitié de deux fois plus
d'\omterm{def:b1:enzymes:enzyme}{enzyme} est toujours active)~; avec I elle double aussi (la vitesse est
proportionnelle à l'\omterm{def:b1:enzymes:enzyme}{enzyme} à tout $[S]$)~: doubler l'\omterm{def:b1:enzymes:enzyme}{enzyme} ne permet
jamais de distinguer les \omterm{def:b1:enzymes:inhibitors}{inhibiteurs}.
\textbf{18.} Vitesse à \qty{25}{\celsius}~: $V_{\max}[S]/(K_m + [S])$
avec $V_{\max} = k_{\text{cat}}[E]$~: par litre, $1670\times 10^{-8} =
\qty{1.67e-5}{mol/L/s}$~; $\times 0.5/2.5 = \qty{3.3e-6}{mol/L/s}$~; à
\qty{37}{\celsius}, $\times 2^{1.2} = 2.3$~: \qty{7.7e-6}{mol/L/s}~; dans
\qty{1e-12}{L}~: \qty{7.7e-18}{mol/s} $= \num{4.6e6}$ molécules par
seconde.
\textbf{19.} Aucune augmentation nécessaire~: la vitesse (\num{4.6e6})
dépasse la demande (\num{3e6})~; s'il fallait l'augmenter, la \omterm{def:b1:cell-unit-of-life:cell}{cellule}
pourrait élever le \omterm{def:b1:enzymes:enzyme}{substrat}, ou phosphoryler l'\omterm{def:b1:enzymes:enzyme}{enzyme} pour abaisser son
$K_m$, ou fabriquer plus d'\omterm{def:b1:enzymes:enzyme}{enzyme}.
\textbf{20.} $K_m^{\text{app}} = 2(1 + 0.4/0.2) = \qty{6}{mmol/L}$~:
vitesse $\times 0.5/6.5$ au lieu de $0.5/2.5$~: \qty{38}{\%} de la
précédente, \num{1.8e6} par seconde --- le produit freine sa propre
synthèse en dessous de la demande, de sorte qu'il est consommé et que
son taux baisse, ce qui libère l'\omterm{def:b1:enzymes:enzyme}{enzyme}~: un approvisionnement qui
s'ajuste tout seul.
\textbf{21.} Avant~: $0.5/2.5 = 0.20$ de $V_{\max}$~; après~: $0.5/0.9 =
0.56$~: presque le triple.
\textbf{22.} Quasi nulle~: à pH 5 l'histidine du \omterm{def:b1:enzymes:enzyme}{site actif} est protonée
et ne peut plus agir en base, et les résidus acides du site sont
neutralisés~; l'état d'ionisation de l'\omterm{def:b1:enzymes:enzyme}{enzyme}, et peut-être son
repliement, sont faussés. Elle est faite pour le \omterm{def:b1:eukaryotic-cell:organelle}{cytosol}, non pour le
\omterm{def:b1:eukaryotic-cell:endomembrane}{lysosome}.
\textbf{23.} Le $K_m$ change peu (la fixation tient surtout au reste du
site)~; le $k_{\text{cat}}$ s'effondre de plusieurs ordres de grandeur,
puisque l'histidine était la base catalytique qui activait l'eau ou la
sérine.
\textbf{24.} Au voisinage de $K_m$, la vitesse répond aux variations de
\omterm{def:b1:enzymes:enzyme}{substrat} (pente voisine de $V_{\max}/2K_m$)~; très au-dessus, l'\omterm{def:b1:enzymes:enzyme}{enzyme}
est saturée et insensible, et l'excès de \omterm{def:b1:enzymes:enzyme}{substrat} serait une charge
osmotique et chimique. Travailler près de $K_m$ rend la voie
commandable par l'approvisionnement.
\textbf{25.} $K_m = \qty{2.0}{mmol/L}$, $V_{\max} = \qty{1.0}{\micro mol/min}$
(\omterm{def:b1:enzymes:enzyme}{enzyme} à \qty{10}{nmol/L}), $k_{\text{cat}} = \qty{1670}{s^{-1}}$~; I
\omterm{def:b1:enzymes:inhibitors}{compétitif}, $K_i = \qty{1.0}{mmol/L}$~; J non \omterm{def:b1:enzymes:inhibitors}{compétitif}, $K_i =
\qty{1.0}{mmol/L}$.
\end{solution}
